\section{מושגים בסיסיים}
לצורך הבנת העבודה, להלן הסבר על המושגים המרכזיים בתחומי האגרונומיה, חיזוי סדרות עיתיות ולמידת מכונה שבהם נעשה שימוש בפרויקט.

\subsection{מושגים אגרונומיים}

\begin{itemize}
\item \textbf{\color{ruppinblue}מוליכות חשמלית \LR{(EC -- Electrical Conductivity)}:} מדד לכמות המלחים המומסים בתמיסת המצע או הקרקע. בחקלאות, EC משמש כאינדיקציה לרמת המליחות ולריכוז הדשן הזמין לצמח. ערכים גבוהים מדי עלולים לגרום לעקה אוסמוטית ולהקשות על קליטת מים, בעוד שערכים נמוכים מדי עשויים להעיד על מחסור בהזנה. המדד נמדד ביחידות mS/cm.

\item \textbf{\color{ruppinblue}חומציות \LR{(pH)}:} מדד לרמת החומציות או הבסיסיות של התמיסה. ערך ה-pH משפיע על מסיסות וזמינות יסודות ההזנה לצמח, ובפרט על זמינות מיקרו-אלמנטים כגון ברזל, מנגן ואבץ. סטייה מטווח ה-pH הרצוי עלולה לגרום למחסורי הזנה או לרעילות, גם כאשר הדשן סופק בכמות מספקת.

\item \textbf{\color{ruppinblue}התאדות-דיות \LR{(ET0 -- Evapotranspiration)}:} מדד המייצג את פוטנציאל איבוד המים מהמערכת כתוצאה מהתאדות ומהדיות של הצמח. ET0 מושפע בעיקר מקרינה, טמפרטורה, לחות ורוח, והוא משמש להערכת דרישת המים של הגידול ולתכנון השקיה.

\item \textbf{\color{ruppinblue}בית שורשים \LR{(Rootzone)}:} נפח המצע או הקרקע שבו מתפתחת מערכת השורשים של הצמח. באזור זה מתרחשים תהליכי קליטת מים ודשן, הצטברות מלחים ושינויים בחומציות. לכן, תנאי בית השורשים משפיעים ישירות על בריאות הצמח, קצב הצימוח והיבול.

\item \textbf{\color{ruppinblue}מצע מנותק:} שיטת גידול שבה הצמח גדל במצע שאינו הקרקע הטבעית, כגון קוקוס, פרלייט או מצעים מלאכותיים אחרים. במצע מנותק נפח בית השורשים מוגבל וקיבולת הבאפר נמוכה יחסית, ולכן שינויים בהשקיה, בדישון או במליחות עשויים להשפיע במהירות על הצמח.
\end{itemize}

\subsection{מושגים בלמידת מכונה}

\begin{itemize}
\item \textbf{\color{ruppinblue}מודל רגרסיה \LR{(Regression Model)}:} מודל חיזוי שמטרתו להעריך ערך מספרי רציף. בפרויקט זה נעשה שימוש במודלי רגרסיה לצורך חיזוי ערכי pH ו-EC, וכן לצורך חיזוי משתני מיקרו-אקלים בתוך החממה.

\item \textbf{\color{ruppinblue}חיישן תוכנה \LR{(Soft Sensor)}:} מודל חיזוי המשמש להערכת משתנה שקשה או יקר למדוד באופן רציף, על בסיס משתנים אחרים שקל יותר למדוד. בפרויקט זה, ה־\SoftSensor{} נועד להעריך את ערכי ה-pH וה-EC בבית השורשים בין מועדי המדידה הידניים, באמצעות נתוני אקלים, השקיה, דישון, מצב הצמח ומדידות עבר.

\item \textbf{\color{ruppinblue}\GB{}:} משפחת אלגוריתמים המבוססת על שילוב הדרגתי של עצי החלטה רבים. כל עץ חדש נבנה במטרה לתקן את השגיאות של העצים הקודמים, וכך מתקבל מודל חיזוי חזק המסוגל ללמוד קשרים מורכבים ולא ליניאריים בין משתנים.

\item \textbf{\color{ruppinblue}\LR{XGBoost}:} מימוש מתקדם ויעיל של \GB{} המבוסס על עצי החלטה. האלגוריתם מתאים במיוחד לנתונים טבלאיים, מאפשר למידה של קשרים לא-ליניאריים, ומתמודד היטב עם שילובים מורכבים בין משתנים. בפרויקט זה שימש XGBoost כרכיב המרכזי במודל החיזוי של בית השורשים.

\item \textbf{\color{ruppinblue}\LR{LightGBM}:} אלגוריתם נוסף ממשפחת \GB{} המתאפיין במהירות אימון גבוהה וביעילות חישובית. בפרויקט זה שימש LightGBM בעיקר בשלבי חיזוי משתני מיקרו-אקלים בתוך החממה, על בסיס נתוני אקלים חיצוני.

\item \textbf{\color{ruppinblue}מודל ליניארי חסין \LR{(Robust Linear Model)}:} מודל רגרסיה ליניארי שנועד להיות פחות רגיש לערכים חריגים ולרעש בנתונים. בפרויקט זה שולב רכיב ליניארי חסין כחלק מהמודל ההיברידי, במטרה לשפר את תגובת החיזוי במצבים של אירועים תפעוליים בעלי השפעה גבוהה, כגון שינוי חד במליחות או מעבר בין מגמות.

\item \textbf{\color{ruppinblue}מודל היברידי:} מודל המשלב יותר מגישת חיזוי אחת. בפרויקט זה המודל ההיברידי משלב בין \GB{}, הלומד דפוסים מורכבים מתוך נתוני העבר, לבין רכיב ליניארי חסין, המסייע בהתמודדות עם מגמות ואירועים חריגים. שילוב זה מאפשר למודל לבצע גם אינטרפולציה בתוך דפוסים מוכרים וגם אקסטרפולציה זהירה במצבים פחות שכיחים.

\item \textbf{\color{ruppinblue}אינטרפולציה:} חיזוי עבור מצב הדומה למצבים שכבר הופיעו בנתוני האימון. לדוגמה, חיזוי תגובת בית השורשים ביום בעל תנאי אקלים והשקיה הדומים לימים שנצפו בעבר.

\item \textbf{\color{ruppinblue}אקסטרפולציה:} חיזוי עבור מצב החורג חלקית מהדפוסים שנצפו בנתוני האימון. בפרויקט זה מושג זה רלוונטי בעיקר לאירועים תפעוליים חריגים או לשינויים חדים במליחות ובחומציות, שבהם נדרש מהמודל להעריך מגמת שינוי גם מעבר לדוגמאות השכיחות בנתונים.

\item \textbf{\color{ruppinblue}\WFV{}:} שיטת אימות המתאימה לסדרות עיתיות. במקום לחלק את הנתונים באופן אקראי לאימון ולבדיקה, המודל מאומן על נתוני העבר ונבדק על נקודות עתידיות. לאחר מכן נקודת הזמן מתקדמת, והמודל נבחן שוב. שיטה זו מדמה תרחיש תפעולי שבו בכל שלב זמינים רק נתוני העבר.

\item \textbf{\color{ruppinblue}\LR{Holdout}:} שיטת בדיקה שבה חלק מהנתונים נשמר מחוץ לתהליך האימון ומשמש להערכת ביצועי המודל לאחר האימון. בפרויקט זה נעשה שימוש גם בבדיקת Holdout על מרווחים מדולגים, כדי לבחון האם המודל מסוגל לחזות נקודות שלא שימשו ללמידה.

\item \textbf{\color{ruppinblue}הנדסת פיצ'רים \LR{(Feature Engineering)}:} תהליך יצירת משתנים חדשים מתוך הנתונים הגולמיים במטרה לשפר את יכולת הלמידה של המודל. דוגמאות מפרויקט זה הן זמן שעבר מאז ההשקיה האחרונה, סכום הקרינה המצטבר, כמות הדשן בפרק זמן נתון, משך הזמן בין מדידת הבסיס לנקודת החיזוי, ומדדי היסטוריה של בית השורשים.

\item \textbf{\color{ruppinblue}\LR{MAE -- Mean Absolute Error}:} מדד שגיאה המחושב כממוצע הערכים המוחלטים של ההפרשים בין הערכים האמיתיים לערכים החזויים. ככל שה-MAE נמוך יותר, כך החיזוי מדויק יותר. היתרון המרכזי של MAE הוא שהוא נמדד באותן יחידות של המשתנה החזוי, למשל יחידות pH או mS/cm עבור EC.

\item \textbf{\color{ruppinblue}\LR{RMSE -- Root Mean Squared Error}:} מדד שגיאה המחושב כשורש ממוצע ריבועי השגיאות. RMSE נותן משקל גבוה יותר לשגיאות גדולות, ולכן הוא מתאים לזיהוי מצבים שבהם המודל מבצע טעויות חריגות.

\item \textbf{\color{ruppinblue}\Rtwo{} (מקדם ההסבר):} מדד המבטא איזה חלק מהשונות בנתוני היעד מוסבר על ידי המודל. ערך הקרוב ל-1 מעיד על התאמה טובה יותר של המודל לנתונים, בעוד שערך נמוך מעיד כי המודל מתקשה להסביר את השונות.
\end{itemize}
